\documentclass[11pt,a4paper]{article}
\usepackage[utf8]{inputenc}
\usepackage[T1]{fontenc}
\usepackage[spanish]{babel}
\usepackage{lmodern}
\usepackage[margin=2.3cm]{geometry}
\usepackage{xcolor,listings,booktabs,longtable,hyperref}
\usepackage{fancyhdr}
\hypersetup{colorlinks=true,linkcolor=blue,urlcolor=blue}
\lstset{basicstyle=\ttfamily\small,breaklines=true,frame=single,columns=fullflexible,showstringspaces=false,inputencoding=utf8,literate={á}{{\'a}}1 {é}{{\'e}}1 {í}{{\'i}}1 {ó}{{\'o}}1 {ú}{{\'u}}1 {ñ}{{\~n}}1 {Á}{{\'A}}1 {É}{{\'E}}1 {Í}{{\'I}}1 {Ó}{{\'O}}1 {Ú}{{\'U}}1 {Ñ}{{\~N}}1}
\pagestyle{fancy}
\fancyhf{}
\lhead{IFCD107 -- Plataforma de trabajo}
\rhead{API y Docker}
\cfoot{\thepage}
\setlength{\headheight}{14pt}
\title{Plataforma de trabajo del curso\\\large API REST, Docker y servicios de apoyo}
\author{IFCD107 -- Especialista en Inteligencia Artificial}
\date{Octubre de 2026}
\begin{document}
\maketitle

\begin{abstract}
Este documento describe el entorno práctico del proyecto de regresión polinómica. La aplicación Flask expone una API HTTP para ejecutar modelos de aprendizaje automático; Docker empaqueta la aplicación y permite iniciar servicios auxiliares como Ollama, n8n y Flowise. Se explican los componentes, el intercambio de datos, la puesta en marcha, las pruebas y la gestión cotidiana del entorno.
\end{abstract}

\section{Objetivo de la plataforma}
El propósito es practicar dos habilidades complementarias: consumir y diseñar servicios mediante una API, y ejecutar aplicaciones en un entorno reproducible mediante contenedores. En vez de llamar directamente a las funciones Python desde cada herramienta, el cliente envía una petición HTTP a Flask. La API valida la entrada, ejecuta el algoritmo y devuelve una respuesta JSON con métricas, identificador de ejecución y, cuando aplica, rutas de los artefactos generados.

Docker permite preparar los componentes con versiones y configuración declaradas. Cada servicio se ejecuta en su contenedor, con sus dependencias aisladas; Docker Compose coordina su arranque, las redes internas, los puertos y los volúmenes persistentes.

\section{Componentes del proyecto}
\begin{longtable}{p{0.24\linewidth}p{0.68\linewidth}}
\toprule
Componente & Función en las prácticas\\\midrule
API Flask (\texttt{app.py}) & Publica rutas HTTP para salud, autenticación y entrenamientos. Delega la lógica de modelos en \texttt{api\_functions.py} y \texttt{regression/}.\\
Dockerfile & Construye la imagen de la API sobre Python 3.11, instala \texttt{requirements.txt} y ejecuta Gunicorn en el puerto 5000.\\
Docker Compose & Declara los servicios y su configuración. El Compose raíz incluye la API, Ollama, n8n y Flowise, además del Compose oficial de Dify cuando su repositorio está preparado.\\
Ollama & Sirve modelos de lenguaje locales en la red de contenedores. La configuración descarga por defecto \texttt{llama3:8b} y \texttt{mistral:7b}.\\
n8n & Entorno visual para construir flujos de automatización que pueden llamar a APIs.\\
Flowise & Interfaz visual para crear flujos con modelos y componentes de IA.\\
Dify & Plataforma de aplicaciones LLM incorporada desde su Compose oficial; requiere preparar su repositorio y entorno.\\
Postman & Cliente para enviar solicitudes HTTP, adjuntar CSV y revisar respuestas sin escribir un cliente propio.\\
\bottomrule
\end{longtable}

La API también se puede ejecutar directamente con Python, sin Docker. Esta opción es útil para desarrollar y depurar Flask; Docker resulta adecuado para practicar el despliegue de servicios relacionados. La colección de Postman está en \texttt{postman/Regresion\_Polinomica.postman\_collection.json} y los ejemplos de datos en \texttt{ejemplos/}.

\section{Modelo mental: imagen, contenedor y servicio}
Una \textbf{imagen} es la plantilla inmutable que contiene el programa y sus dependencias. Un \textbf{contenedor} es una instancia en ejecución de esa imagen. Un \textbf{servicio} de Compose describe cómo crear y ejecutar uno o más contenedores: imagen o construcción, variables, puertos, volúmenes y dependencias.

Los servicios de Compose se comunican entre sí por la red interna usando el nombre del servicio como dirección. Por ejemplo, la API usa \texttt{http://ollama:11434} para alcanzar Ollama. Desde el navegador del equipo, en cambio, se usa un puerto publicado como \texttt{http://localhost:5000}. El puerto interno y el publicado se conectan mediante una regla del Compose.

\section{Recorrido de una petición a la API}
\begin{enumerate}
\item Postman envía una solicitud HTTP a la dirección local de la API.
\item Flask identifica la ruta y el método, valida los datos y, para las rutas protegidas, comprueba el token Bearer.
\item La función correspondiente prepara los datos, separa entrenamiento y prueba, ajusta el modelo y calcula predicciones y métricas.
\item El proyecto guarda los resultados bajo \texttt{runs/<referencia>/}, incluyendo según el endpoint registros, CSV, gráficas o modelos.
\item Flask devuelve un estado HTTP y un cuerpo JSON. Postman muestra ambos para inspección.
\end{enumerate}

Una API es un contrato entre cliente y servidor. El método indica la operación; la ruta identifica el recurso o acción; los encabezados transportan metadatos como el tipo de contenido y la autorización; el cuerpo contiene la entrada. JSON se utiliza para datos estructurados y \texttt{multipart/form-data} para formularios que incluyen ficheros.

\section{Rutas representativas}
\begin{longtable}{p{0.14\linewidth}p{0.43\linewidth}p{0.36\linewidth}}
\toprule Método & Ruta & Uso\\\midrule
GET & \texttt{/api/v1/health} & Comprobar que la API responde; no requiere token.\\
POST & \texttt{/api/v1/auth/token} & Enviar credenciales y obtener un token temporal.\\
GET & \texttt{/api/v1/regresion/parametros} & Consultar el catálogo de parámetros disponibles.\\
POST & \texttt{/api/v1/regresion/polinomica} & Ajustar el modelo a datos JSON.\\
POST & \texttt{/api/v1/regresion/polinomica/csv} & Ajustar el modelo a un CSV con columnas \texttt{x,y}.\\
POST & \texttt{/api/v1/regresion/polinomica/randomizedsearch} & Buscar hiperparámetros con validación cruzada.\\
POST & \texttt{/api/v1/random-forest/csv} & Entrenar Random Forest desde un CSV tabular.\\
\bottomrule
\end{longtable}
Las rutas de entrenamiento requieren \texttt{Authorization: Bearer <token>}. La emisión del token y la comprobación de salud son públicas. El código de estado ayuda a reconocer el resultado: 2xx indica éxito, 400 entrada inválida, 401 falta de autorización, 413 exceso de tamaño y 5xx un fallo del servidor o del entrenamiento.

\subsection{Ejemplo de solicitud JSON}
\begin{lstlisting}
POST http://localhost:5000/api/v1/regresion/polinomica
Authorization: Bearer <token>
Content-Type: application/json

{
  "reference": "prueba_curso",
  "degree": 2,
  "data": [
    {"x": 1, "y": 3.1},
    {"x": 2, "y": 5.2},
    {"x": 3, "y": 9.0},
    {"x": 4, "y": 15.1},
    {"x": 5, "y": 22.8}
  ]
}
\end{lstlisting}
Los campos exactos y límites se deben consultar en la guía del proyecto y en la colección Postman. Para CSV, Postman construye automáticamente el formulario multipart al adjuntar el archivo; no hay que fijar manualmente el encabezado boundary.

\section{Puesta en marcha con Docker}
El entorno descrito usa Docker Desktop con Docker Compose v2.20 o posterior. La primera ejecución descarga imágenes y modelos; Ollama puede necesitar varios gigabytes de almacenamiento, además de memoria y tiempo de CPU/GPU. Se necesita conexión a Internet para esas descargas iniciales.

El Compose raíz requiere los archivos de configuración locales. Si no existen, crear el entorno a partir de las plantillas; si ya existen, conservarlos y completar las variables que falten. El repositorio oficial de Dify y su archivo de entorno son necesarios si se va a levantar Dify:
\begin{lstlisting}
if (!(Test-Path .env)) { Copy-Item .env.example .env }
git clone https://github.com/langgenius/dify.git dify
if (!(Test-Path dify/docker/.env)) {
  Copy-Item dify/docker/.env.example dify/docker/.env
}
\end{lstlisting}
Antes de iniciar, configurar credenciales y secretos aleatorios propios en los archivos \texttt{.env}; no publicar esos archivos ni copiar secretos a documentación compartida. Para Dify, la guía principal incluye la configuración de perfiles, sandbox y puertos. Si solo se quiere ejecutar la API localmente, se puede omitir este conjunto de servicios y usar el procedimiento Python de \texttt{README.md}.

Para construir e iniciar los servicios:
\begin{lstlisting}
docker compose up -d --build
docker compose ps
docker compose logs -f regression-api
\end{lstlisting}
La opción \texttt{-d} deja los contenedores en segundo plano; \texttt{--build} reconstruye la imagen de la API. La descarga de modelos se realiza mediante el servicio de inicialización de Ollama y puede prolongar el primer arranque. Los servicios principales se publican localmente en:
\begin{center}
\begin{tabular}{ll}\toprule
Servicio & Dirección\\\midrule
API de regresión & \url{http://localhost:5000}\\
n8n & \url{http://localhost:5678}\\
Flowise & \url{http://localhost:3001}\\
Dify (si se configuró) & \url{http://localhost:8080}\\
Ollama & \url{http://localhost:11434}\\\bottomrule
\end{tabular}
\end{center}
Las interfaces locales están ligadas a la máquina anfitriona. Para comprobar la API, abrir \url{http://localhost:5000/api/v1/health} o ejecutar \texttt{Invoke-RestMethod http://localhost:5000/api/v1/health} en PowerShell.

\section{Práctica guiada en Postman}
\begin{enumerate}
\item Importar la colección de Postman del proyecto.
\item Configurar \texttt{baseUrl} como \texttt{http://localhost:5000}, el usuario \texttt{admin} y la contraseña tomada del entorno local autorizado.
\item Ejecutar primero \textbf{Obtener token}. La colección guarda el token para las siguientes peticiones.
\item Ejecutar \textbf{Health} y comprobar el estado HTTP y el JSON.
\item Probar el ajuste polinómico JSON y, después, el ajuste CSV con \texttt{ejemplos/datos\_ejemplo\_xy.csv}.
\item Inspeccionar métricas, \texttt{run\_id} y archivos bajo \texttt{runs/}. Cambiar el grado o la referencia y repetir.
\item Probar una entrada inválida para observar cómo la API devuelve un error HTTP legible.
\end{enumerate}

\section{Datos, persistencia y limpieza}
El servicio de la API monta un volumen Docker llamado \texttt{regression\_runs} sobre \texttt{/app/runs}; así, los artefactos persisten aunque se reemplace el contenedor. Ollama guarda sus modelos en \texttt{./ollama}; n8n y Flowise tienen volúmenes propios para sus datos y configuración. Eliminar un contenedor no equivale a borrar el volumen.

Para detener el conjunto manteniendo los volúmenes:
\begin{lstlisting}
docker compose down
\end{lstlisting}
Para iniciarlo otra vez, ejecutar \texttt{docker compose up -d}. Evitar \texttt{docker compose down -v} si se quieren conservar modelos, flujos, configuraciones y resultados. Para revisar un servicio concreto, usar \texttt{docker compose logs --tail 100 nombre-del-servicio}.

\section{Buenas prácticas de trabajo}
\begin{itemize}
\item Mantener claves, contraseñas y tokens en los archivos locales de entorno; compartir solo plantillas sin secretos.
\item Usar referencias sencillas para agrupar ejecuciones y no incluir datos personales en ellas.
\item Empezar con los CSV sintéticos de \texttt{ejemplos/}; están preparados para practicar y no representan registros reales.
\item Verificar el estado del servicio y los logs antes de atribuir un error al algoritmo.
\item Distinguir la dirección usada por el equipo (\texttt{localhost}) de la dirección usada entre contenedores (nombre de servicio, por ejemplo \texttt{ollama}).
\item Recordar que un modelo local no envía la inferencia a un proveedor externo, pero sí utiliza recursos del equipo.
\end{itemize}

\section{Resultados de aprendizaje}
Al trabajar con esta plataforma, el alumnado practica la lectura de contratos HTTP, el uso de métodos y códigos de estado, el intercambio de JSON y CSV, la autenticación Bearer, el diagnóstico mediante logs y la ejecución de servicios con Docker Compose. En paralelo, puede observar cómo una API integra el entrenamiento de modelos y convierte sus resultados en respuestas y artefactos consultables.

\section{Documentación del proyecto}
Este texto sirve como presentación de la plataforma. Para detalles de implementación y contratos completos, consultar \texttt{documentacion/guia\_proyecto.tex}; para instrucciones actualizadas de arranque, configuración y solución de problemas, consultar \texttt{README.md}. Las solicitudes reproducibles están en la colección de Postman.

\end{document}
